Table~\ref{tab:group_diff} lists the features with notable adoption
differences across the three groups (spread $> 10$ percentage points).
Features with similar adoption rates across groups are omitted from the table
and discussed below.

\begin{table}[ht]
  \centering
  \caption{Features with divergent adoption rates across groups
           (spread $> 10$ pp).  The highest value per row is set in
           \textbf{bold}.}
  \label{tab:group_diff}
  \begin{tabular}{llrrr}
    \toprule
    \textbf{Category} & \textbf{Feature} &
      \textbf{Internal} & \textbf{MINT} & \textbf{Community} \\
    \midrule
    \multirow{5}{*}{\textcolor{catPresentation}{Presentation}}
      & Narrator       & \textbf{86.6} &  9.4 & 84.0 \\
      & TTS fragments  & \textbf{36.1} &  0.1 & 12.6 \\
      & Animations     & \textbf{66.2} &  0.6 & 25.9 \\
      & Anim.\ blocks  & \textbf{36.8} &  0.3 & 11.3 \\
      & Galleries      & \textbf{17.3} &  0.0 &  9.5 \\
    \midrule
    \multirow{4}{*}{\textcolor{catInteraction}{Interaction}}
      & Any quiz       & 22.7 & \textbf{78.3} & 29.6 \\
      & Single choice  & 11.9 &  0.2 & \textbf{14.1} \\
      & Multiple choice& 11.3 &  0.8 & \textbf{16.3} \\
      & Text quiz      &  9.7 & \textbf{70.0} &  6.2 \\
    \midrule
    \multirow{3}{*}{\textcolor{catReuse}{Reuse}}
      & Macros (any)   & 69.0 & \textbf{99.7} & 54.4 \\
      & Imports        & 57.4 & \textbf{98.3} & 34.3 \\
      & Ext.\ scripts  & 16.2 & \textbf{96.3} & 30.3 \\
    \midrule
    \multirow{5}{*}{\textcolor{catEmbedding}{Embedding}}
      & Exec.\ code    & \textbf{11.7} &  0.3 &  2.7 \\
      & WebApps        & \textbf{13.0} &  0.7 &  6.9 \\
      & Script tags    & \textbf{33.3} &  8.2 & 25.8 \\
      & ASCII diagrams & \textbf{38.5} &  2.3 &  7.1 \\
      & Math           & 30.7 & \textbf{95.2} & 21.0 \\
      & HTML embeds    & 18.6 &  5.9 & \textbf{21.5} \\
    \bottomrule
  \end{tabular}
\end{table}

Three patterns emerge from these divergences:

\paragraph{MINT-the-GAP's template-driven workflow.}
MINT-the-GAP dominates the Reuse category (macros~\SI{99.7}{\percent},
imports~\SI{98.3}{\percent}, external scripts~\SI{96.3}{\percent})
and relies heavily on text quizzes (\SI{70.0}{\percent}) and
math notation (\SI{95.2}{\percent}).
This profile reflects a standardised production pipeline for
school-level STEM content, where a fixed template provides presentation
and reuse infrastructure while quizzes serve as the primary interactive
element.
Notably, presentation features like narrator (\SI{9.4}{\percent}) and
animations (\SI{0.6}{\percent}) are nearly absent, suggesting that the
template workflow bypasses them entirely.

\paragraph{Internal as feature showcase.}
The Internal group leads in presentation features (narrator~\SI{86.6}{\percent},
animations~\SI{66.2}{\percent}) and embedding features
(ASCII diagrams~\SI{38.5}{\percent}, script tags~\SI{33.3}{\percent}).
This is consistent with the developers' role as feature demonstrators and
early adopters of the full language surface.

\paragraph{Community as selective adopters.}
The Community group generally falls between the other two, but with
notable exceptions: single-choice (\SI{14.1}{\percent}) and
multiple-choice quizzes (\SI{16.3}{\percent}) are \emph{higher} than
in Internal, suggesting that community authors prefer familiar quiz
types over text-input quizzes.
HTML embeds (\SI{21.5}{\percent}) also peak in the Community group.

\paragraph{Convergent features.}
Several features show similar adoption across all three groups:
matrix quizzes (\SIrange{0.3}{2.4}{\percent}), surveys
(\SIrange{0.0}{4.8}{\percent}), task lists (\SIrange{0.0}{8.7}{\percent}),
and code projects (\SIrange{0.0}{1.3}{\percent}).
These consistently low rates indicate features that remain underutilised
regardless of author background, pointing to adoption barriers that
transcend group-specific workflows.
